\section{Related literature}
\label{sec:related}

\aidraft{%
\textbf{Regime and jump-model allocation.}
The direct targets are the statistical-jump-model allocation papers: \citet{nystrup2020eswa,nystrup2021}
on the estimator, and \citet{shu2024} (single-asset downside overlay) and \citet{shu2025dynamic}
(multi-asset, expected-return-conditioned optimization), with the model-predictive-control extension
\citep{jmmpc2025}.
Our quarrel is not with the estimator --- which we adopt and validate --- but with the baselines these
papers use to establish economic value: buy-and-hold, minimum-variance, and equal-weight, none
exposure-matched or reactive.%
}

\aidraft{%
\textbf{The deflation genre.}
\citet{cederburg2020} show that volatility-managed portfolios fail out of sample across a large factor
cross-section; our result is adjacent --- a regime-conditioning strategy that dies against the same
reactive incumbents --- and inherits the same warning about a factor-timing chapter.%
}

\aidraft{%
\textbf{Reactive incumbents.}
The bars we hold the label to are the volatility-timing and volatility-managed literatures
\citep{fko2001,moreiramuir2017}, with \citet{barroso2015} the robust exception (momentum
volatility-scaling), and the time-series-momentum construction of \citet{moskowitz2012} underlying the
trend proxy of Section~\ref{sec:race4}.
The momentum-crash dynamic of \citet{danielmoskowitz2016} appears in the descriptive anatomy.%
}
\needcite{DeMiguel et al.\ 2024 if used; Kritzman--Page--Turkington 2012 (Regime Shifts).}
